% Options for packages loaded elsewhere
\PassOptionsToPackage{unicode}{hyperref}
\PassOptionsToPackage{hyphens}{url}
\PassOptionsToPackage{dvipsnames,svgnames,x11names}{xcolor}
\documentclass[
  10pt,
  fontset=fandol]{ctexart}
\usepackage{xcolor}
\usepackage[margin=16mm]{geometry}
\usepackage{amsmath,amssymb}
\setcounter{secnumdepth}{-\maxdimen} % remove section numbering
\usepackage{iftex}
\ifPDFTeX
  \usepackage[T1]{fontenc}
  \usepackage[utf8]{inputenc}
  \usepackage{textcomp} % provide euro and other symbols
\else % if luatex or xetex
  \usepackage{unicode-math} % this also loads fontspec
  \defaultfontfeatures{Scale=MatchLowercase}
  \defaultfontfeatures[\rmfamily]{Ligatures=TeX,Scale=1}
\fi
\usepackage{lmodern}
\ifPDFTeX\else
  % xetex/luatex font selection
\fi
% Use upquote if available, for straight quotes in verbatim environments
\IfFileExists{upquote.sty}{\usepackage{upquote}}{}
\IfFileExists{microtype.sty}{% use microtype if available
  \usepackage[]{microtype}
  \UseMicrotypeSet[protrusion]{basicmath} % disable protrusion for tt fonts
}{}
\makeatletter
\@ifundefined{KOMAClassName}{% if non-KOMA class
  \IfFileExists{parskip.sty}{%
    \usepackage{parskip}
  }{% else
    \setlength{\parindent}{0pt}
    \setlength{\parskip}{6pt plus 2pt minus 1pt}}
}{% if KOMA class
  \KOMAoptions{parskip=half}}
\makeatother
\setlength{\emergencystretch}{3em} % prevent overfull lines
\providecommand{\tightlist}{%
  \setlength{\itemsep}{0pt}\setlength{\parskip}{0pt}}
\usepackage{amsmath,amssymb,mathtools,needspace}
\xeCJKDeclareCharClass{CJK}{"2032,"0400->"04FF,"2160->"216B,"2460->"2473,"25A0,"25C7,"25CB,"25CF}
\IfFontExistsTF{Arial Unicode MS}{\xeCJKsetup{AutoFallBack=true}\setCJKfallbackfamilyfont{\CJKrmdefault}{Arial Unicode MS}}{}
\setcounter{secnumdepth}{0}
\setlength{\emergencystretch}{2em}
\allowdisplaybreaks
\usepackage{bookmark}
\IfFileExists{xurl.sty}{\usepackage{xurl}}{} % add URL line breaks if available
\urlstyle{same}
\hypersetup{
  pdftitle={小结},
  colorlinks=true,
  linkcolor={Maroon},
  filecolor={Maroon},
  citecolor={Blue},
  urlcolor={Blue},
  pdfcreator={LaTeX via pandoc}}

\title{小结}
\author{}
\date{}

\begin{document}
\maketitle

\subsection{小结}\label{ux5c0fux7ed3}

本章研究求解常微分方程的差分方法。构造差分格式主要有两条途径：基于数值积分的构造方法和基于 Taylor 展开的构造方法。后一种方法更灵活，也更具有一般性。Taylor 展开方法还有一个优点，它在构造差分公式的同时可以得到关于截断误差的估计。

但是，直接用 Taylor 展开导出的 Taylor 级数法（见 7.3.1 节）不便于实际应用。基于 Taylor 展开构造出的四阶 Runge-Kutta 方法（见 7.3.4 节）则是电子计算机上的常用算法。四阶 Runge-Kutta 方法的优点是精度高，程序简单，计算过程稳定，并且易于调节步长。

四阶 Runge-Kutta 方法也有不足之处，它要求函数 \(f\) 具有较高的光滑性。如果 \(f\) 的光滑性差，那么它的精度可能还不如 Euler 格式（见 7.2.1 节）或改进的 Euler 格式（见 7.2.4 节）。

四阶 Runge-Kutta 方法的另一个缺点是计算量比较大，需要耗费较多的机器时间（每一步需四次计算函数 \(f\) 的值）。相比之下，Hamming 方法（见 7.5.7 节）可以节省计算量（每一步只需两次计算函数 \(f\) 的值）。但 Hamming 方法是一种四步法，它不是自开始的，需要借助于四阶 Runge-Kutta 方法（或四阶 Taylor 级数法）提供开始值。

就差分方法来说，常微分方程与偏微分方程有着紧密的联系，读者可参看偏微分方程的数值解法的有关教材。

\begin{center}\rule{0.5\linewidth}{0.5pt}\end{center}

\subsection{相关原页校注（agent 补充）}\label{ux76f8ux5173ux539fux9875ux6821ux6ce8agent-ux8865ux5145}

以下校注来自本节覆盖的原页，可能也涉及同页相邻小节；原文照录与校正说明分开保留。

\subsubsection{原书 PDF 页 153 的校注}\label{ux539fux4e66-pdf-ux9875-153-ux7684ux6821ux6ce8}

\href{../pages/pdf-153.md}{完整原页转录} · \href{../../source-images/pdf-153.jpeg}{原页图像}

校注记录：ch05b-E007。

\begin{quote}
校注（agent 补充，ch05b-E007）：小结中的五处节号按原印刷保留为 7.3.1、7.3.4、7.2.1、7.2.4、7.5.7；这些节号与本版章节编排不符。对应主题在本书章节索引中分别为 5.3.1（Taylor 级数法）、5.3.5（四阶 Runge-Kutta 方法）、5.2.1（Euler 格式）、5.2.4（改进的 Euler 格式）、5.5.7（Hamming 格式）。其中本章 PDF 140 的正文也明确将四阶 Runge-Kutta 格式指向 5.3.5，PDF 144 的实际标题为 5.5.7 Hamming 格式。
\end{quote}

\subsection{本节覆盖原页的校注记录}\label{ux672cux8282ux8986ux76d6ux539fux9875ux7684ux6821ux6ce8ux8bb0ux5f55}

下列记录按原页关联，含同页相邻小节的校注；计算时同时读取上文校注和完整原页。

\begin{itemize}
\tightlist
\item
  \texttt{ch05b-E007}：原书 PDF 页 153
\end{itemize}

\end{document}
